% !TeX root = ../main.tex

\section{Revisión de la literatura}
\label{sec:revision-literatura}

\subsection{Una relación establecida y sus condiciones}
\label{subsec:literatura-relacion}

La bibliografía seleccionada permite distinguir dos preguntas: cómo influye la dispersión del servicio bajo condiciones constantes y qué intervenciones pueden mejorar un sistema concreto. La primera tiene una respuesta analítica para M/G/1 FCFS; la segunda exige identificar qué cambia realmente en el proceso. Esta revisión acotada organiza los aportes según esa distinción y sus límites, sin pretender una búsqueda sistemática ni una cobertura exhaustiva del campo.

La relación central no constituye una hipótesis teórica nueva. La ecuación~\ref{eq:pollaczek-khinchine}, desarrollada en el marco conceptual, establece que, a tasa de llegada y media de servicio fijas, una mayor varianza aumenta la espera media en cola, siempre que se cumplan los supuestos y la finitud de momentos indicados allí \cite{virtamoMG1}. El apéndice~1 de Joustra y colaboradores recupera esta relación para explicar el efecto de la variabilidad del servicio y la comparación entre servicio constante y exponencial \cite{joustra2012radioterapia}. Su utilidad para esta monografía es ofrecer una referencia común con la cual interpretar los demás trabajos, no una regla independiente de las condiciones del modelo.

El preprint de Lauber Bonomo, Pal y Reuveni también parte de M/G/1 FCFS y expresa el número medio de trabajos en el sistema mediante la utilización y el coeficiente de variación del servicio \cite{lauberbonomo2021resetting}. Es importante distinguir sus indicadores: en la sección~II, el tiempo denominado \emph{waiting time} comprende desde la llegada hasta el final del servicio. Corresponde, por tanto, a $T$ en nuestra notación, no a $W_q$. Con media de servicio fija, la diferencia entre ambos promedios es constante; cuando una intervención modifica esa media, la distinción vuelve a ser decisiva.

\subsection{Aplicaciones: variabilidad de capacidad y métodos complementarios}
\label{subsec:literatura-aplicaciones}

En procesos de varias etapas, hablar de variabilidad requiere precisar su origen. Joustra y colaboradores distinguen la dispersión de las duraciones individuales de servicio de las fluctuaciones de la capacidad disponible. En su estudio de radioterapia, la primera se relaciona con la espera en sala, mientras que la segunda afecta el acceso a la primera consulta y el flujo hacia preparación y tratamiento \cite{joustra2012radioterapia}. Esta separación evita interpretar cualquier reducción de demoras como evidencia de una disminución de la varianza del tiempo de atención.

El estudio combina análisis de datos del proceso, teoría de colas y simulación de eventos discretos. La teoría proporciona mecanismos explicativos; la simulación representa las interacciones entre etapas, localiza restricciones de capacidad y compara escenarios. Los autores contrastan el modelo con datos de la práctica para las medias y los percentiles considerados, y emplean intervalos de confianza al evaluar el cumplimiento de metas. En los escenarios estudiados, una capacidad de primeras consultas más estable permite mejorar el acceso y las demoras posteriores \cite{joustra2012radioterapia}. La evidencia cuantitativa procede del modelo de simulación contrastado con el proceso, no de un ensayo operativo que aísle únicamente la dispersión del servicio.

El límite de transferencia es concreto: se trata de un proceso con citas, recursos múltiples y etapas interdependientes, y varios escenarios incluyen un acelerador adicional. El propio artículo aclara que no interviene sobre la duración de las consultas, cuyas fluctuaciones eran pequeñas. Así, su aporte no consiste en comprobar experimentalmente nuestra comparación M/G/1 a media fija, sino en mostrar por qué conviene combinar explicación analítica y representación dinámica cuando la organización del servicio excede una cola aislada \cite{joustra2012radioterapia}. Para la demostración futura de esta monografía, la combinación tendrá un propósito más limitado: ilustrar una identidad conocida con parámetros hipotéticos, sin afirmar validación de un servicio real.

\subsection{Intervenciones que modifican el servicio efectivo}
\label{subsec:literatura-reinicios}

Otra forma de actuar sobre la congestión consiste en modificar el mecanismo de servicio. Lauber Bonomo y colaboradores analizan reinicios cuando la aleatoriedad es intrínseca al servidor: al reiniciar un trabajo, se obtiene una nueva duración aleatoria independiente, aunque el tiempo ya invertido se pierde. No equivale a adelantar otro cliente ni a eliminar trabajos pendientes. Si, por el contrario, cada trabajo conserva un requerimiento fijo y el reinicio solo obliga a repetirlo, el mecanismo puede desperdiciar tiempo sin aportar esa posibilidad de renovación \cite{lauberbonomo2021resetting}.

La estrategia matemática consiste en caracterizar el tiempo total de servicio efectivo, incluidos los intentos fallidos, y volver a aplicar el análisis M/G/1 con sus nuevos momentos. El preprint estudia reinicios a tasa constante y a intervalos fijos, y desarrolla ejemplos lognormales y de Pareto mediante fórmulas y evaluación numérica. Bajo las condiciones examinadas, la elección adecuada del reinicio puede reducir la media del servicio y su variabilidad relativa, con mejoras en el número medio de trabajos \cite{lauberbonomo2021resetting}. Son resultados analíticos y ejemplos numéricos del trabajo citado, no simulaciones realizadas en esta monografía.

Esta intervención no mantiene necesariamente constante $m$: también puede alterar $\rho=\lambda m$. En consecuencia, una mejora no debe atribuirse exclusivamente a la reducción de la varianza. Tampoco contradice que el servicio determinista minimice la espera media entre distribuciones de igual media: después del reinicio, la media efectiva puede ser diferente. El aporte del estudio es mostrar que dispersión y velocidad media pueden cambiar conjuntamente; nuestra comparación las separará deliberadamente \cite{lauberbonomo2021resetting}.

\subsection{Límites al cambiar estructura, disciplina o indicador}
\label{subsec:literatura-limites}

La estructura de la red también restringe las generalizaciones. En el resumen y el comienzo de la introducción de Bondi y Whitt se presentan principios sobre variabilidad para una red cerrada, de una clase de trabajos, con estaciones de un servidor y disciplina FCFS. Allí se anuncia que la longitud media en una estación que constituye un cuello de botella puede disminuir al aumentar la variabilidad del servicio en esa estación \cite{bondiwhitt1986variabilidad}. Esta referencia se utiliza únicamente como contraste de alcance, a partir de esa primera página consultada; no se evalúan aquí sus demostraciones, ejemplos numéricos ni procedimientos de aproximación. Una red cerrada no conserva el proceso externo de llegadas Poisson de nuestra cola abierta aislada, por lo que el principio anunciado no puede trasladarse directamente a ella.

El cambio de disciplina y de indicador ofrece un contraste más explícito. Koole, Nuyens y Righter estudian el máximo número de clientes en el sistema durante un período ocupado en M$^X$/G/1, con llegadas Poisson por lotes y variantes LCFS con interrupción: se atiende primero a la llegada más reciente. Sus resultados comparan distribuciones de ese máximo, no esperas medias FCFS. Para sistemas estables, igual tasa de llegada e igual distribución de tamaño del lote, obtienen máximos estocásticamente menores cuando los servicios están ordenados como más variables en sentido convexo \cite{koole2005maximos}.

El orden convexo exige comparar las esperanzas de todas las funciones convexas para las cuales existan; implica igualdad de medias y, cuando las varianzas son finitas, ordena también sus varianzas. Sin embargo, una varianza mayor por sí sola no establece ese orden. Por ello, los teoremas no autorizan a predecir el máximo a partir de cualquier comparación de varianzas. El manuscrito aporta caracterizaciones recursivas, demostraciones para las disciplinas preemptivas estudiadas e ilustraciones numéricas; reserva otras disciplinas para conjeturas \cite{koole2005maximos}. No corresponde presentar estas últimas como resultados demostrados ni interpretar un máximo menor como una mejora universal de todos los indicadores.

\subsection{Síntesis y fundamento de la comparación controlada}
\label{subsec:literatura-sintesis}

En conjunto, los trabajos no ofrecen conclusiones incompatibles: responden preguntas diferentes. La relación de referencia concierne a la espera media de una cola abierta FCFS; la aplicación hospitalaria modifica disponibilidad de capacidad; los reinicios transforman el servicio efectivo; los otros contrastes cambian la estructura de red o la disciplina y el indicador. Su lectura conjunta muestra que una afirmación sobre variabilidad debe identificar siempre qué se mantiene fijo, qué mecanismo se modifica y qué desempeño se mide.

La monografía se justifica, entonces, como una síntesis explicativa y una comparación controlada, no como la resolución de un vacío teórico. Mantener $\lambda$, $m$ y FCFS constantes permitirá aislar el efecto del segundo momento sobre $\mathrm{E}[W_q]$. La futura demostración deberá contrastar estimaciones con la referencia analítica y considerar su incertidumbre, sin equiparar una ejecución finita con la media estacionaria. La metodología precisará ese diseño; aquí queda fundamentado por qué conviene preservar estos controles y no extender sus conclusiones a redes, reinicios, máximos de cola o disciplinas no estudiadas.
